\documentclass[11pt,a4paper]{article}
% Même mise en forme que le polycopié de prérequis (boîtes, exercices, corrigés).
% ===========================================================================
%  Polycopié de prérequis — Comptabilité prudentielle (M2 alternance)
%  Préambule partagé
%  F. Borel-Mathurin — Université Paris 1 Panthéon-Sorbonne
% ===========================================================================

\usepackage[utf8]{inputenc}
\usepackage[T1]{fontenc}
\IfFileExists{french.ldf}{\usepackage[french]{babel}}{}
\IfFileExists{lmodern.sty}{\usepackage{lmodern}}{}
\usepackage[a4paper,margin=2.4cm]{geometry}

\usepackage{amsmath,amssymb}
\usepackage{booktabs}
\usepackage{array}
\usepackage{enumitem}
\usepackage{xcolor}
\usepackage{tcolorbox}
\tcbuselibrary{skins,breakable}
\usepackage{tikz}
\usetikzlibrary{arrows.meta,positioning,calc}
\usepackage{fancyhdr}
\usepackage[hidelinks]{hyperref}
\usepackage{makeidx}
\makeindex
\IfFileExists{eurosym.sty}{\usepackage{eurosym}}{\providecommand{\euro}{\texteuro}}

% ---------------------------------------------------------------- couleurs
% À remplacer par les valeurs exactes de la charte si besoin.
\definecolor{bleuactu}{RGB}{0,84,150}
\definecolor{vertactu}{RGB}{0,132,100}
\definecolor{rougeactu}{RGB}{175,30,45}
\definecolor{orangeactu}{RGB}{230,140,0}
\definecolor{grisactu}{RGB}{110,110,110}

% ------------------------------------------------- version étudiant/corrigé
% \correctiontrue  -> corrigés inclus
% \correctionfalse -> version étudiant
\newif\ifcorrection
\correctiontrue

% ------------------------------------------------------------------ boîtes
\newtcolorbox{notion}[1]{%
  breakable, enhanced, colback=bleuactu!4, colframe=bleuactu,
  boxrule=0.4pt, left=8pt, right=8pt, top=6pt, bottom=6pt,
  fonttitle=\bfseries\sffamily, coltitle=white, title={#1}}

\newtcolorbox{generalisation}[1]{%
  breakable, enhanced, colback=orangeactu!5, colframe=orangeactu!85!black,
  boxrule=0.4pt, left=8pt, right=8pt, top=6pt, bottom=6pt,
  fonttitle=\bfseries\sffamily, coltitle=white, title={#1}}

\newtcolorbox{pont}[1]{%
  breakable, enhanced, colback=rougeactu!4, colframe=rougeactu,
  boxrule=0.4pt, left=8pt, right=8pt, top=6pt, bottom=6pt,
  fonttitle=\bfseries\sffamily, coltitle=white, title={#1}}

\newtcolorbox{refreg}{%
  breakable, enhanced, colback=grisactu!6, colframe=grisactu!60,
  boxrule=0.4pt, left=8pt, right=8pt, top=6pt, bottom=6pt,
  fontupper=\small}

% --------------------------------------------------------------- exercices
\newcounter{exo}[section]
\renewcommand{\theexo}{\thesection.\arabic{exo}}

\newenvironment{exercice}[1][]{%
  \refstepcounter{exo}%
  \par\medskip\noindent
  \begin{tcolorbox}[breakable, enhanced, colback=white, colframe=vertactu,
     boxrule=0.4pt, left=8pt, right=8pt, top=6pt, bottom=6pt,
     fonttitle=\bfseries\sffamily, coltitle=white,
     title={Exercice \theexo\ifx\relax#1\relax\else\ --- #1\fi}]}%
  {\end{tcolorbox}\par\medskip}

\newenvironment{corrige}{%
  \ifcorrection
    \par\smallskip\noindent
    \begin{tcolorbox}[breakable, enhanced, colback=vertactu!4,
       colframe=vertactu!50, boxrule=0.4pt, left=8pt, right=8pt,
       top=6pt, bottom=6pt, fontupper=\small,
       title={\itshape Corrigé}, fonttitle=\small\itshape,
       coltitle=vertactu!80!black, colbacktitle=vertactu!10]%
  \else
    \setbox0\vbox\bgroup
  \fi}%
  {\ifcorrection\end{tcolorbox}\par\smallskip\else\egroup\fi}

% ------------------------------------------------------------- échéanciers
\tikzset{
  timeaxis/.style={-{Latex[length=2mm]}, thick, grisactu},
  timenode/.style={circle, fill=bleuactu, inner sep=1.1pt},
  flux/.style={-{Latex[length=2mm]}, thick, vertactu},
  fluxneg/.style={-{Latex[length=2mm]}, thick, rougeactu},
  etiq/.style={font=\scriptsize, inner sep=1pt}
}

% ------------------------------------------------------------------ titres
\usepackage{titlesec}
\titleformat{\section}{\Large\bfseries\sffamily\color{bleuactu}}{\thesection}{0.7em}{}
\titleformat{\subsection}{\large\bfseries\sffamily\color{bleuactu!80!black}}{\thesubsection}{0.6em}{}

\pagestyle{fancy}
\fancyhf{}
\fancyhead[L]{\small\sffamily\color{grisactu}Prérequis --- Comptabilité prudentielle}
\fancyhead[R]{\small\sffamily\color{grisactu}\thepage}
\renewcommand{\headrulewidth}{0.2pt}

\newcommand{\notation}[2]{\item[{$#1$}] #2}  % accolades : #1 peut contenir des crochets


% Version étudiant par défaut ; `make pdf` produit aussi td-corrige.pdf en définissant \Corrige.
\ifdefined\Corrige\correctiontrue\else\correctionfalse\fi

\fancyhead[L]{\small\sffamily\color{grisactu}Séance 1 --- Comptabilité prudentielle}
\newcommand{\keur}{\,k\euro}
\newcommand{\meur}{\,M\euro}

\title{\vspace{-1.5cm}\sffamily\bfseries\color{bleuactu}
       Cinq cas pour entrer dans le cours\\[2mm]
       \large\mdseries\color{black}Pourquoi, comment, quand, à quoi ça sert, et ses limites
       \ifcorrection\\[1mm]\color{vertactu}Éléments de mise en commun\fi}
\author{\sffamily Fabrice Borel-Mathurin}
\date{\sffamily\small Master SAFIR --- 2026-2027}

\begin{document}
\maketitle
\thispagestyle{empty}

\begin{notion}{Comment travailler (1 h 30, puis 1 h de mise en commun)}
Par groupes de trois ou quatre, chaque groupe prend un cas. Vingt minutes pour répondre aux
questions avec ce que vous savez déjà et un peu de bon sens : il n'y a pas de calcul à faire. Puis
chaque groupe présente son cas en cinq minutes aux autres. Chaque cas annonce une séance du cours :
notez les questions qui restent ouvertes, nous y reviendrons.
\end{notion}

\section{Les cas}

\begin{exercice}[A. Pourquoi ? Equitable Life, Royaume-Uni]
Equitable Life est une mutuelle d'assurance vie britannique fondée en 1762. Des années 1950 à la
fin des années 1980, elle vend des contrats de retraite assortis d'une option : à l'échéance,
l'assuré peut convertir son épargne en rente à un taux garanti. À la vente, les taux de marché sont
élevés : l'option ne vaut presque rien. Dans les années 1990, les taux baissent et l'option devient
très coûteuse. L'entreprise tente de réduire la participation aux bénéfices des assurés qui
l'exercent ; la justice le lui interdit en 2000, et elle cesse la même année de vendre des contrats.
\begin{enumerate}[nosep]
  \item Pourquoi le coût de cette option n'apparaissait-il pas dans les comptes au moment de la vente ?
  \item Que faudrait-il mesurer, et à quelle fréquence, pour le voir venir ?
  \item Dans une mutuelle, qui supporte finalement ce coût ?
\end{enumerate}
\end{exercice}
\begin{corrige}
\begin{enumerate}[nosep]
  \item Des provisions calculées avec des hypothèses fixées à la souscription, au coût historique,
        ne donnent aucune valeur à une option « hors de la monnaie » : elle n'apparaît qu'une fois
        exercée.
  \item La valeur \emph{actuelle} de l'option, dans les conditions de marché du jour : c'est ce que
        fait la meilleure estimation de Solvabilité II, calculée sur de nombreux scénarios de taux
        (séances 5 à 7), et recalculée au moins une fois par an.
  \item Les autres assurés, dont la participation aux bénéfices baisse. D'où l'objectif principal
        du contrôle : protéger les assurés (directive 2009/138/CE, art.~27).
\end{enumerate}
\end{corrige}

\begin{exercice}[B. Comment ? L'Assureur Bêta]
Un assureur vie fictif affiche 50\,000\keur{} de capitaux propres dans ses comptes sociaux. Passé
au bilan prudentiel, il ressemble à ceci : ses placements valent 26\,000\keur{} de moins que leur
valeur comptable ; ses provisions techniques, recalculées comme une meilleure estimation plus une
marge de risque, valent 10\,000\keur{} de plus ; un impôt différé actif de 9\,000\keur{} apparaît.
Son SCR vaut 70\,000\keur.
\begin{enumerate}[nosep]
  \item Que valent ses fonds propres prudentiels (excédent d'actif sur passif) ?
  \item D'où peut venir chacun des trois écarts ?
  \item À première vue, son SCR est-il couvert ? Que faudrait-il savoir de plus pour conclure ?
\end{enumerate}
\end{exercice}
\begin{corrige}
\begin{enumerate}[nosep]
  \item $50\,000 - 26\,000 - 10\,000 + 9\,000 = 23\,000$\keur.
  \item Moins-values latentes (taux, écarts de crédit, actions) ; provisions comptables moins
        élevées que la valeur de transfert des engagements, surtout du fait de la marge de risque ;
        impôt différé actif sur ces moins-values et ce surcroît de provisions, retenu s'il est
        recouvrable.
  \item Non : 23\,000 contre 70\,000. Pour conclure, il faut savoir quels fonds propres comptent
        (niveaux, limites d'éligibilité) et si une dette subordonnée vient s'ajouter. C'est le TD
        de la séance 2 : le ratio y ressort à 70\,\%.
\end{enumerate}
\end{corrige}

\begin{exercice}[C. Quand ? Le cycle des taux, 2012--2022]
Un assureur vie a vendu des contrats d'épargne rachetables à tout moment, avec un taux minimum
garanti de 2 à 3\,\%. Ses placements sont surtout des obligations longues. De 2012 à 2021, les taux
de marché baissent jusqu'à devenir négatifs sur certaines échéances ; en 2022, ils remontent
brutalement.
\begin{enumerate}[nosep]
  \item Quand les taux baissent, que deviennent la valeur des obligations, la meilleure estimation
        des contrats et le ratio de couverture ?
  \item Quand les taux remontent vite, quel risque nouveau apparaît ?
  \item À quel rythme l'assureur et l'ACPR voient-ils ces effets dans les chiffres prudentiels ?
\end{enumerate}
\end{exercice}
\begin{corrige}
\begin{enumerate}[nosep]
  \item Les obligations prennent de la valeur, mais la meilleure estimation augmente aussi :
        les flux futurs sont moins actualisés et la garantie de taux coûte plus cher. Si les
        engagements sont plus longs que les placements, le passif monte plus que l'actif et le
        ratio baisse. Tout dépend de l'écart de duration (séances 6 et 7).
  \item Les rachats : les assurés partent vers des placements mieux rémunérés, et l'assureur
        doit parfois vendre des obligations en moins-value pour les payer (séance 5).
  \item En continu en interne ; le MCR est recalculé chaque trimestre, le SCR au moins une fois
        par an et dès que le profil de risque change sensiblement (directive, art.~102 et 129) ;
        les états remis au superviseur sont trimestriels et annuels.
\end{enumerate}
\end{corrige}

\begin{exercice}[D. À quoi ça sert ? L'Assureur Alpha]
L'Assureur Alpha a 120\meur{} de fonds propres éligibles et un SCR de 100\meur{} : ratio de 120\,\%.
Un choc de marché lui coûte 20\meur{} de fonds propres. Son conseil d'administration devait voter
le mois suivant un dividende de 10\meur.
\begin{enumerate}[nosep]
  \item Quel est son ratio après le choc ? Le conseil peut-il voter le dividende ?
  \item Citez trois leviers pour remonter le ratio. Lesquels agissent sur les fonds propres,
        lesquels sur le SCR ?
  \item Quel ratio viseriez-vous en temps normal, et qui le décide ?
\end{enumerate}
\end{exercice}
\begin{corrige}
\begin{enumerate}[nosep]
  \item 100\,\% en supposant le SCR inchangé : l'entreprise est au seuil. Un dividende prévisible
        se déduit des fonds propres : le ratio tomberait à 90\,\% et le SCR ne serait plus couvert.
        Le dividende n'est pas raisonnable.
  \item Sur les fonds propres : augmentation de capital, émission de dette subordonnée, mise en
        réserve du résultat. Sur le SCR : réassurance, réduction de l'exposition aux actions,
        couvertures, produits avec moins de garanties.
  \item Nettement plus de 100\,\%, pour encaisser un choc sans intervention du superviseur. La
        cible relève de l'appétit au risque fixé par le conseil, éclairé par l'ORSA (séance 9).
\end{enumerate}
\end{corrige}

\begin{exercice}[E. Les limites : le printemps 2020]
En mars 2020, la pandémie de Covid-19 fait chuter les marchés d'actions ; les écarts de crédit
s'élargissent ; les taux restent très bas. Deux mécanismes prévus par le cadre amortissent la baisse
des ratios : l'ajustement symétrique, qui réduit le choc actions du SCR après une baisse des
marchés, et la correction pour volatilité, qui relève le taux d'actualisation des provisions quand
les écarts de crédit s'élargissent. En avril, l'EIOPA appelle les assureurs à suspendre leurs
dividendes.
\begin{enumerate}[nosep]
  \item Pourquoi un bilan en valeur de marché peut-il amplifier une crise s'il ne prévoit rien ?
  \item Que corrigent ces deux mécanismes ? Quel en est le prix ?
  \item Le SCR est calibré sur des données passées. Après un choc record, mesure-t-il encore le bon risque ?
\end{enumerate}
\end{exercice}
\begin{corrige}
\begin{enumerate}[nosep]
  \item Les ratios baissent tous en même temps ; pour les remonter, les assureurs vendent les
        actifs risqués, ce qui fait encore baisser les prix : c'est la procyclicité.
  \item L'ajustement symétrique allège le choc actions après une baisse (et l'alourdit après une
        hausse) ; la correction pour volatilité neutralise une partie de l'élargissement
        temporaire des écarts de crédit. Le prix : un ratio moins fidèle à la valeur de marché
        du jour, plus difficile à comparer, et des paramètres à justifier (séances 6 à 8).
  \item Pas forcément : les paramètres sont estimés sur un historique qui ne contenait pas un tel
        choc, et la question de les réévaluer après coup se pose. C'est un sujet de recherche, et
        l'une des raisons de l'ORSA, qui demande à chaque entreprise d'évaluer elle-même son
        besoin de capital.
\end{enumerate}
\end{corrige}

\begin{generalisation}{Après la séance}
Notez, pour chaque cas, la question qui vous paraît encore ouverte. Chaque séance du cours en
reprendra une ; la carte du cours, sur la dernière diapositive de la séance, dit laquelle.
\end{generalisation}

\end{document}
